\subsection{THE INVARIANTS THEOREM}

Let g be a Lie algebra and \(\rho\) a representation of g on a vector space M. For every class \(\delta\) of simple representations of g let \(M_{\delta}\) be the isotypical component of M of species \(\delta\) . The subspace \(M_{0}\) of invariant elements of M is just \(M_{\delta_{0}}\) , where \(\delta_{0}\) denotes the class of the zero representation of g on a space of dimension 1.

Lemma 4. Let \(\rho, \sigma, \tau\) be representations of \(\mathfrak{g}\) on vector spaces \(\mathbf{M}, \mathbf{N}, \mathbf{P}\) . Suppose that we are given a K-bilinear mapping \((m, n) \mapsto m \cdot n\) of \(\mathbf{M} \times \mathbf{N}\) into \(\mathbf{P}\) such that

\[
(\rho (x) m). n + m. (\sigma (x) n) = \tau (x) (m. n)
\]

for all \(m \in M\) , \(n \in N\) , \(x \in g\) .

\begin{enumerate}
\def\labelenumi{(\alph{enumi})}
\item
  If \(m_0 \in \mathbf{M}_0\) , the mapping \(n \mapsto m_0\) . \(n\) is a \(\mathfrak{g}\) -module homomorphism.
\item
  If \(n \in N_{\delta}\) , then \(m_{0}.n \in P_{\delta}\) .
\item
  If \(\mathbf{M}\) is a (not necessarily associative) algebra and the \(\rho(x)\) are derivations of \(\mathbf{M}\) , \(\mathbf{M}_0\) is a subalgebra of \(\mathbf{M}\) and each \(\mathbf{M}_{\delta}\) is a right and left \(\mathbf{M}_0\) -module.
\end{enumerate}

For \(m_{0} \in M_{0}, n \in N\) and \(x \in g,\)

\[
\tau (x) \left(m _ {0}. n\right) = m _ {0}. (\sigma (x) n),
\]

whence (a). Assertion (b) follows from (a) (Algebra, Chapter VIII, § 3, no. 4, Proposition 10). If \(\mathbf{N} = \mathbf{P} = \mathbf{M}\) and \(\sigma = \tau = \rho\) , assertion (b) gives assertion (c) as a special case.

Lemma 5. Suppose further that \(\sigma\) and \(\tau\) are semi-simple and hence N (resp. P) is the direct sum of the \(N_{\delta}\) (resp. \(P_{\delta}\) ). For all \(n \in N\) (resp. \(p \in P\) ), let \(n^{\sharp}\) (resp. \(p^{\sharp}\) ) be its component in \(N_0\) (resp. \(P_0\) ). Let \(m_0 \in M_0\) . Then for all \(n \in N\) , \((m_0.n)^{\sharp} = m_0.n^{\sharp}\) .

By linearity it suffices to consider the case where \(n \in \mathrm{N}_{\delta}\) . If \(\delta \neq \delta_0\) , \(n^{\sharp} = 0\) and \(m_0.n \in \mathrm{P}_{\delta}\) (Lemma 4), hence \((m_0.n)^{\sharp} = 0 = m_0.n^{\sharp}\) . If \(\delta = \delta_0\) , \(n^{\sharp} = n\) and \(m_0.n \in \mathrm{P}_0\) (Lemma 4), hence \((m_0.n)^{\sharp} = m_0.n = m_0.n^{\sharp}\) .

THEOREM 6. Let \(\mathfrak{g}\) be a Lie algebra, V a semi-simple \(\mathfrak{g}\) -module of finite dimension over K, S the symmetric algebra of V and \(x_{\mathrm{S}}\) the derivation of S which extends \(x_{\mathrm{V}}\) (so that \(x \mapsto x_{\mathrm{S}}\) is a representation of \(\mathfrak{g}\) on S).

\begin{enumerate}
\def\labelenumi{(\alph{enumi})}
\item
  The algebra \(\mathbf{S}_0\) of invariants of \(\mathbf{S}\) is generated by a finite number of elements.
\item
  For every class \(\delta\) of simple representations of \(\mathfrak{g}\) of finite dimension over \(\mathbf{K}\) , let \(\mathrm{S}_{\delta}\) be the isotypical component of \(\mathrm{S}\) of species \(\delta\) . Then \(\mathrm{S}_{\delta}\) is a finitely generated \(\mathrm{S}_{0}\) -module.
\end{enumerate}

Let \(\overline{\mathbf{S}}\subset \mathbf{S}\) be the ideal of elements of S with no constant term. Let I be the ideal of S generated by \(\mathrm{S}_0\cap \overline{\mathrm{S}}\) and let \((s_1,s_2,\ldots ,s_p)\) be a finite system of generators of the ideal I (Commutative Algebra, Chapter III, § 3). It can be assumed that the \(s_t\) belong to \(\mathbf{S}_0 \cap \bar{\mathbf{S}}\) and are homogeneous (the \(x_{\mathrm{S}}\) preserve degrees and hence each \(\mathbf{S}_0\) is a graded submodule). Let \(\mathbf{S}_1\) be the subalgebra of \(\mathbf{S}\) generated by 1 and the \(s_t\) . Then \(\mathbf{S}_1 \subset \mathbf{S}_0\) . We show that \(\mathbf{S}_1 = \mathbf{S}_0\) . For this we prove that every homogeneous element \(s\) of \(\mathbf{S}_0\) is in \(\mathbf{S}_1\) , arguing by induction on the degree \(n\) of \(s\) . If \(n = 0\) , our assertion is obvious. Suppose therefore \(n > 0\) and that our assertion has been proved when the degree of s is \(< n\) . As \(s \in I\) , \(s = \sum_{i=1}^{p} s_i s_i'\) , where the \(s_i'\) are elements of S which can be assumed to be homogeneous, with \(\deg(s_i') = \deg(s) - \deg(s_i) < n\) . Lemma 5 can be applied, as the g-module S is semi-simple (no. 5, Corollary 2 to Theorem 4); in the notation of this lemma,

\[
s = s ^ {\sharp} = \sum_ {i = 1} ^ {p} (s _ {i} s _ {i} ^ {\prime}) ^ {\sharp} = \sum_ {i = 1} ^ {p} s _ {i} {s _ {i} ^ {\prime}} ^ {\sharp}.
\]

The \(s_i^\sharp\) are elements of \(S_0\) which are homogeneous and of degree \(< n\) (since each \(S_\delta\) is a graded submodule). Hence they are in \(S_1\) by the induction hypothesis. Hence \(s \in S_1\) , which proves (a).

We now consider a simple representation of g of class \(\delta\) on a finite-dimensional space M. Let \(L = \mathcal{L}_{\mathbb{K}}(M, S)\) . For all \(s \in S\) and all \(f \in L\) , let sf be the element of L defined by \((sf)(m) = s.f(m)\) ( \(m \in M\) ); an S-module structure is thus defined on L; as M is finite-dimensional over K, clearly L is a finitely generated S-module and hence a Noetherian S-module since the ring S is Noetherian. On the other hand, L has a canonical g-module structure. For every integer \(n \geqslant 0\) let \(S^{n}\) be the set of homogeneous elements of S of degree n; then the g-module \(\mathcal{L}_{\mathbb{K}}(M, S^{n})\) is semi-simple (no. 5, Corollary 3 to Theorem 4) and hence the g-module L is semi-simple. Moreover, for \(s \in S\) , \(f \in L\) , \(x \in g\) and \(m \in M\) ,

\[
\begin{array}{r l} (x _ {\mathrm{L}} (s f)) (m) & = x _ {\mathrm{S}} ((s f) (m)) - (s f) (x _ {\mathrm{M}} m) \\ & = x _ {\mathrm{S}} (s. f (m)) - s. f (x _ {\mathrm{M}} m) \\ & = (x _ {\mathrm{S}} s). f (m) + s. (x _ {\mathrm{S}} f (m)) - s. f (x _ {\mathrm{M}} m) \\ & = ((x _ {\mathrm{S}} s) f) (m) + (s (x _ {\mathrm{L}} f)) (m) \end{array}
\]

and hence \(x_{\mathrm{L}}(sf) = (x_{\mathrm{S}}s)f + s(x_{\mathrm{L}}f)\) . We can therefore apply Lemma 5.

The subset \(\mathbf{L}_0\) of invariant elements of \(\mathbf{L}\) is just the set of homomorphisms of the \(g\) -module \(\mathbf{M}\) into the \(g\) -module \(\mathbf{S}\) . Hence, if \(\phi\) denotes the canonical homomorphism of \(\mathrm{M} \otimes_{\mathrm{K}} \mathrm{L}\) onto \(\mathrm{S}\) , \(\phi(\mathrm{M} \otimes_{\mathrm{K}} \mathrm{L}_0) = \mathrm{S}_{\delta}\) . As \(\phi\) is obviously an S-module homomorphism, it suffices to show that \(\mathbf{L}_0\) is a finitely generated \(\mathrm{S}_0\) -module. Let \(\mathbf{J}\) be the sub-S-module of \(\mathbf{L}\) generated by \(\mathbf{L}_0\) . Since \(\mathbf{L}\) is a Noetherian S-module, there exists a finite sequence \((f_1, \ldots, f_q)\) of elements of \(\mathbf{L}_0\) generating the S-module \(\mathbf{J}\) . Let \(\mathbf{L}_1\) be the \(\mathrm{S}_0\) -module generated by \(f_{1},\ldots ,f_{q}\) . Then \(\mathrm{L}_1\subset \mathrm{L}_0\) . On the other hand, if \(f\in \mathrm{L}_0\) , then \(f = \sum_{i = 1}^{q}s_{i},f_{i}\) with \(s_i\in S\) for all \(i\) and hence by Lemma 5 whose notation we adopt:

\[
f = f ^ {\sharp} = \left(\sum_ {i = 1} ^ {q} s _ {i} f _ {i}\right) ^ {\sharp} = \sum_ {i = 1} ^ {q} s _ {i} ^ {\sharp} f _ {i} \in L _ {1}.
\]

Hence \(L_{0} = L_{1}\) , so that \(L_{0}\) is a finite generated \(S_{0}\) -module.

